\documentclass[11pt]{article}
\usepackage[margin=1in]{geometry}
\usepackage{amsmath,amssymb,amsthm}
\usepackage{booktabs,array}
\usepackage{graphicx}
\usepackage{microtype}
\usepackage{xcolor}
\usepackage[hidelinks]{hyperref}

\newtheorem{theorem}{Theorem}
\newtheorem{lemma}{Lemma}
\newcommand{\SAP}{\mathrm{SAP}}
\newcommand{\Sol}{\operatorname{Sol}}
\newcommand{\im}{\operatorname{im}}

\title{SAP Is Equivalent to Lunarity for Coordinatewise-Injective
$4\times4$ Maps}
\author{Automated mathematics run \texttt{fa\_banach\_001}\\
Candidate full result for expert review}
\date{22 July 2026}

\begin{document}
\maketitle

\begin{center}
\fbox{\parbox{0.92\textwidth}{
\textbf{Status.} Candidate full solution, likely valid.  The converse is a
computer-assisted finite proof with exact integer and rational arithmetic.
All 14,828 non-lunar symmetry classes have exact norm certificates.  Human
review and an author-level novelty check are still required.}}
\end{center}

\section{Source problem}

Han, Qiu, and Wang ask in Problem~7 of \cite{HanQiuWang2024}, on PDF
page~21, whether a map
$\Phi:\{1,2,3,4\}\times\{1,2,3,4\}\to\mathcal L$ whose associated Hankel
system has SAP must be lunar.  In the standing setup of that section, the
associated system is defined for coordinatewise-injective maps.  The decisive
sentence is:
\begin{quote}
Assume that the Hankel system associated to $\Phi$ has SAP.  Does it follow
that $\Phi$ is a lunar map?
\end{quote}

\begin{figure}[ht]
\centering
\includegraphics[width=\textwidth]{figures/open_problem_crop.png}
\caption{Definition, preceding theorem, and Problem~7 in the source paper,
PDF page~21.}
\end{figure}

Write $[4]=\{0,1,2,3\}$.  Coordinatewise injectivity says that every color
appears at most once in each row and each column.  Put
\[
 \Gamma_\ell(a,x)=\mathbf 1_{\{\Phi(a,x)=\ell\}}.
\]
For $a,b\in[4]$, define
\[
 \Sol(a,b)=\{(x,y)\in[4]^2:\Phi(a,x)=\Phi(b,y)\}.
\]
The source paper's equivalent formulation says that $\Phi$ is lunar exactly
when any two sets $\Sol(a,b)$ and $\Sol(c,d)$ are equal or disjoint.

At matrix level $r$, for coefficients $C_\ell\in M_r$, set
\begin{align*}
 A_C&=\sum_\ell \Gamma_\ell\otimes C_\ell,\\
 B_C&=\sum_\ell (\Gamma_\ell\otimes\Gamma_\ell)\otimes C_\ell.
\end{align*}
Tensor flips change neither norm nor the SAP question, so this factor order is
equivalent to the usual coefficient-first convention.  SAP requires
$\|A_C\|=\|B_C\|$ for every $r$ and every coefficient family.

\section{Main result}

\begin{theorem}\label{thm:main}
Let $\Phi:[4]\times[4]\to\mathcal L$ be coordinatewise injective.  The
associated Hankel system has SAP if and only if $\Phi$ is lunar.
\end{theorem}

The implication ``lunar implies SAP'' is Theorem~B of
\cite{HanQiuWang2024}.  The new content is the converse.

\section{Proof intuition}

Regard the $4\times4$ table as a proper edge-coloring of $K_{4,4}$.  A lunar
failure is witnessed by two unequal solution sets that overlap.  The two row
pairs, and independently the two column pairs, have only four possible
directed incidence patterns.  Keeping all sixteen combined patterns gives a
complete normal form without assuming that eight indices are distinct.

An exact restricted-growth recursion then lists every proper coloring and
quotients the non-lunar ones by row permutations, column permutations, color
relabeling, and transposition.  This leaves 14,828 classes.  For every class,
integer scalar or $2\times2$ matrix coefficients make an absorbed support
block strictly larger than the original operator.  Exact Sturm counts put the
two squared norms on opposite sides of a rational number.

\section{Complete directed normal forms}

\begin{lemma}\label{lem:normal}
Every non-lunar coordinatewise-injective $4\times4$ table is equivalent under
row and column permutations to a table satisfying one of sixteen fixed
directed failure patterns.
\end{lemma}

\begin{proof}
Choose unequal overlapping solution sets $\Sol(a,b)$ and $\Sol(c,d)$.  After
interchanging them if necessary, choose
\[
 (x,y)\in\Sol(a,b)\cap\Sol(c,d),\qquad
 (z,w)\in\Sol(a,b)\setminus\Sol(c,d).
\]
Thus
\begin{align}\label{eq:failure}
 \Phi(a,x)&=\Phi(b,y), & \Phi(c,x)&=\Phi(d,y),\notag\\
 \Phi(a,z)&=\Phi(b,w), & \Phi(c,z)&\ne\Phi(d,w).
\end{align}

The ordered row edges $(a,b)$ and $(c,d)$ are distinct non-loops.  If they had
the same initial vertex, the first two equalities and injectivity in column
$y$ would force the same terminal vertex; the solution sets would then be the
same.  The same argument with column $x$ rules out a common terminal vertex.
Consequently, after relabeling rows, the two directed edges have one of the
following forms:
\[
 \begin{array}{c|c}
 \text{type}&(a,b,c,d)\\ \hline
 R&(0,1,1,0)\\
 C_+&(0,1,2,0)\\
 C_-&(0,1,1,2)\\
 D&(0,1,2,3).
 \end{array}
\]
They are respectively reverse, crossed in either orientation, or disjoint.

The points $(x,y)$ and $(z,w)$ are distinct non-loop directed column edges.
Because both lie in $\Sol(a,b)$, equality of their initial vertices forces
equality of their terminal vertices by row injectivity, and conversely.
Hence the same four directed types apply independently to $(x,y,z,w)$.
Substitution of the resulting row and column quadruples into
\eqref{eq:failure} gives the claimed $4\times4=16$ fixed patterns.
\end{proof}

Retaining both crossed orientations is essential: simultaneous reversal flips
both, but does not remove their relative orientation.  The verifier therefore
uses the redundant sixteen-pattern list directly.

\section{Exact classification}

The classifier processes the sixteen cells in row-major order.  A color is
represented by its first-occurrence label.  At each cell the recursion either
uses an earlier color absent from the current row and column, or introduces
the next new label.  Every partition of the cells into matchings has exactly
one such restricted-growth representation.  This proves completeness before
any symmetry quotient.

For each table, the code constructs all sixteen sets $\Sol(a,b)$ as bit masks
and tests lunarity exactly.  For tables satisfying a fixed pattern from
Lemma~\ref{lem:normal}, it minimizes the restricted-growth word over all
$24^2$ row and column permutations, both transpose choices, and canonical
color relabeling.  The exact census is:

\begin{center}
\small
\begin{tabular}{@{}rrrr@{}}
\toprule
colors & canonical proper & canonical non-lunar & non-lunar orbits\\
\midrule
4  & 24        & 0         & 0\\
5  & 4,896     & 4,752     & 16\\
6  & 185,532   & 183,168   & 233\\
7  & 1,569,864 & 1,563,168 & 1,564\\
8  & 4,552,617 & 4,495,344 & 4,238\\
9  & 5,787,840 & 5,504,064 & 5,087\\
10 & 3,717,416 & 3,071,652 & 2,862\\
11 & 1,310,592 & 764,760   & 736\\
12 & 265,110   & 78,072    & 87\\
13 & 31,152    & 2,736     & 5\\
14 & 2,076     & 0         & 0\\
15 & 72        & 0         & 0\\
16 & 1         & 0         & 0\\
\midrule
total & 17,427,192 & 15,667,716 & 14,828\\
\bottomrule
\end{tabular}
\end{center}

For additional auditability, the sixteen fixed-pattern counts are shown
below.  They overlap; their union contains 336,208 restricted-growth tables.

\begin{center}
\begin{tabular}{@{}c rrrr@{}}
\toprule
row type $\backslash$ column type & $R$ & $C_+$ & $C_-$ & $D$\\
\midrule
$R$   & 0      & 0      & 0      & 37,180\\
$C_+$ & 0      & 35,439 & 39,043 & 34,477\\
$C_-$ & 0      & 39,043 & 35,439 & 34,477\\
$D$   & 37,180 & 34,477 & 34,477 & 42,329\\
\bottomrule
\end{tabular}
\end{center}

Transposition is a valid orbit symmetry even for matrix witnesses.  If
$C_\ell$ witnesses a representative, use $C_\ell^{\mathsf T}$ for its
transpose.  Then both $A_C$ and $B_C$ are fully transposed, so their norms are
unchanged.

\section{Exact norm certificates}

The support graph of $B_C$ has row vertices $(a,b)$ and column vertices
$(x,y)$, with a block $C_\ell$ precisely when
$\Phi(a,x)=\Phi(b,y)=\ell$.  Its connected components are orthogonal
row-column blocks.  The diagonal component is a copy of $A_C$.  If any other
component $D_C$ satisfies
\begin{equation}\label{eq:gap}
 \|D_C\|>\|A_C\|,
\end{equation}
then $\|B_C\|>\|A_C\|$ and SAP fails.

For 7,331 orbit representatives, scalar coefficients suffice.  The remaining
7,497 use coefficients $C_\ell\in M_2(\mathbb Z)$.  Candidate coefficients
were found numerically, but a record was stored only after the following exact
gate passed.  For integer matrices $A=A_C$ and $D=D_C$, form
\[
 p_A(t)=\det(tI-A^{\mathsf T}A),\qquad
 p_D(t)=\det(tI-D^{\mathsf T}D).
\]
The record contains a rational $q$ that is not a root of either polynomial.
Exact Sturm counts prove
\[
 \#\{\lambda\in\operatorname{roots}(p_A):\lambda>q\}=0,
 \qquad
 \#\{\lambda\in\operatorname{roots}(p_D):\lambda>q\}\ge1.
\]
Therefore $\|A\|^2<q<\|D\|^2$, proving \eqref{eq:gap} without a
floating-point acceptance step.

For example, orbit~0 has representative
\[
 \begin{pmatrix}
 0&1&2&3\\1&0&3&2\\2&3&0&1\\3&2&1&4
 \end{pmatrix}
\]
and scalar coefficients $(1,1,1,-1,1)$.  The exact separating threshold is
$528077641/100000000$; the source polynomial has no root above it and the
absorbed polynomial has one.  Matrix record~86 illustrates the level-two
case; its complete integer coefficients, component index, threshold, and root
counts are retained in the JSON certificate file rather than rounded for
display.

\section{Proof of Theorem~\ref{thm:main}}

\begin{proof}
If $\Phi$ is lunar, its Hankel system has SAP by Theorem~B of
\cite{HanQiuWang2024}.

Conversely, suppose $\Phi$ is non-lunar.  Lemma~\ref{lem:normal} puts a row
and column permutation of its table into one of the sixteen directed failure
normal forms.  The exact enumeration assigns that table, after color
relabeling and possibly transposition, to one of the 14,828 orbit
representatives.  The certificate attached to that representative supplies
scalar or $2\times2$ integer coefficients satisfying \eqref{eq:gap} exactly.
Row and column permutations, color relabeling, and the coefficient-transpose
rule transfer the strict norm gap back to $\Phi$.  Hence the associated system
does not have SAP.  This proves the contrapositive SAP implies lunar.
\end{proof}

\section{Reproducibility and adversarial review}

Run from the packet's \texttt{code/} directory:
\begin{verbatim}
conda run --no-capture-output -n sandbox python \
  verify_full_problem7.py
\end{verbatim}
The verifier recompiles the classifier, checks all census constants, checks
properness and non-lunarity independently in Python, verifies exact and
disjoint certificate coverage, rebuilds the integer block matrices, and
repeats every Sturm computation.  Its final line is:
\begin{verbatim}
PASS: every non-lunar coordinatewise-injective 4x4 map fails SAP
\end{verbatim}

Three discarded routes are important to the audit trail.  First, certificates
for merge-minimal non-lunar coarsenings do not lift automatically, because
tensoring a merged color creates cross terms.  Second, a lunar failure need
not use four distinct rows and four distinct columns; the earlier
disjoint/disjoint census of 6,301 classes was incomplete.  Third, quotienting
the crossed incidence too aggressively missed a relative-orientation case.
The final sixteen-pattern classifier was built after these corrections and
all final certificates were remapped or regenerated against its output.

\section{Limitations, novelty, and review}

This is a computer-assisted exact proof, not a conceptual classification of
why each non-lunar pattern violates SAP.  It relies on the source theorem for
the forward implication.

A bounded novelty search on 22 July 2026 checked arXiv and general web results
using the exact title, exact Problem~7 wording, ``lunar map'', ``Hankel SAP
converse'', ``prob-hsap'', and the author names.  It found the original
arXiv:2407.01048v1 and a 2024 talk announcement mentioning unspecified
progress, but no later public theorem resolving Problem~7.  Novelty confidence
is medium pending contact with the authors.

The result should be reviewed by an operator-space expert and independently
audited by someone comfortable with finite classification code.  The priority
checks are Lemma~\ref{lem:normal}, transposition with matrix coefficients, and
random replay of certificate records.

\begin{thebibliography}{9}
\bibitem{HanQiuWang2024}
Y.~Han, Y.~Qiu, and Z.~Wang,
\emph{Self-absorption of Hankel systems on monoids -- a seemingly universal
property}, arXiv:2407.01048v1, 2024.
\end{thebibliography}

\end{document}
